\section{Comparative Evaluation: Socio-Financial Agent Dynamics}\label{sec:evaluation}
The preceding section defined the execution loop, budget boundaries, and
deterministic replay mechanisms that advance a simulation through time. To
evaluate how those mechanisms perform in practice, this section examines a
matched empirical benchmark comparing \textit{MiroFish}'s emergent social
discourse \cite{mirofish} with \textit{Boyun}'s stateful financial clearinghouse
on a fourteen-day (336-hour) continuous simulation scenario.

To establish a fair experimental baseline from those execution rules, both
systems were supplied with identical prospective scenario inputs based on the
\textit{MACRON} case across two macroeconomic regimes: \textit{BULL} and
\textit{BEAR}. The evaluation tracks 16 modeled persona agents over 336 hourly
rounds, measuring both qualitative discourse generation and numerical financial
execution. Table~\ref{tab:comparison-benchmark} reports the audited empirical
metrics reconciled across both simulation engines.

\begin{table}[tbp]
\caption{Empirical comparison across 336 simulated hours.}
\label{tab:comparison-benchmark}
\centering\small
\renewcommand{\arraystretch}{1.12}
\begin{tabularx}{\columnwidth}{@{}>{\raggedright\arraybackslash}p{0.36\columnwidth}YY@{}}
\toprule
Dimension & \textit{MiroFish} (BULL / BEAR) & \textit{Boyun} (Live / Stress) \\
\midrule
Simulated horizon & 336h / 336h & 30h--134h / 336h \\
Active personas & 16 / 16 & 16 / 16 \\
Recorded actions & 1,905 / 1,591 & 14,381 / 28,800+ \\
-- Original posts & 537 / 371 & -- \\
-- Quoted debates & 396 / 206 & -- \\
-- Passive waits & 819 / 859 & -- \\
Filled market buys & 0 / 0 & 4,491 / 18,240 \\
Filled market sells & 0 / 0 & 8,808 / 36,480 \\
Rejected executions & 0 / 0 & 1,082 / 4,128 \\
Replay barriers & Uncheckpointed RNG & 20,881 / 41,762 \\
Dormant hours & 86 / 86 (25.6\%) & Continuous \\
Matched policy branches & Unmodeled & 12 / 24 branches \\
Staged exit advantage & Undefined & +\$3,172 to +\$12,625 \\
\bottomrule
\end{tabularx}
\end{table}

Following these reconciled totals, the experimental record reveals two
fundamentally distinct paradigms of agent simulation. In \textit{MiroFish},
persona agents produce organic, contextual narrative text: personas actively
quote opposing market commentators, debate catalyst timing, and express
collective skepticism or euphoria. Over 336 hourly rounds, this generates
537 original posts and 396 quotes under the \textit{BULL} regime, compared to
371 posts and 206 quotes under \textit{BEAR}. However, the environment lacks
an economic clearing mechanism: orders are not settled, token balances and
automated market maker (\textit{AMM}) reserves are unmodeled, and automated
report synthesis timed out during entity-graph memory retrieval.

In contrast to pure discourse generation, \textit{Boyun}'s execution engine
enforces formal \textit{economic settlement} and conservation of capital. Over
its 24 live policy paths, the execution engine recorded 13,299 filled trades
alongside 1,082 rejected executions caused by balance exhaustion or slippage
thresholds. These constraints reflect real market friction that discourse-only
models omit. Furthermore, 20,881 state barriers were verified through
deterministic replay, confirming that identical decision histories produce
identical financial outcomes without re-invoking model inference.

Beyond individual trade execution, the policy comparison mechanisms introduced
in Section~\ref{sec:method-system} require evaluating \textit{counterfactual policy branches}.
\textit{Boyun} evaluated 12 matched policy pairs at identical elapsed
timestamps, comparing a staged exit policy against an inventory-holding pause
policy. Staged liquidation realized between \$3,172.94 and \$12,625.14 more
cash than the pause policy across all twelve matched comparisons. Because pause
retains token inventory, this does not prove higher total wealth; rather, it
provides the decision maker with an auditable trade-off between realized
liquidity and residual inventory risk.

Synthesizing these comparative findings demonstrates that neither qualitative
social emergence nor isolated financial mechanics alone suffices for complex
decision analysis. An expressive multi-agent architecture must couple
\textit{MiroFish}'s emergent narrative discourse with \textit{Boyun}'s
verifiable financial clearinghouse, ensuring that every published message stems
from a bounded information state and every simulated trade adheres to strict
conservation laws.
